\textbf{COLUNI 2021 -- Questão 02}\par\medskip Leia o texto:

\begin{center}
\textit{Sonho Imigrante}

A terra do sonho é distante

e seu nome é Brasil

plantarei a minha vida

debaixo de céu anil.

Minha Itália, Alemanha

Minha Espanha, Portugal

talvez nunca mais eu veja

minha terra natal.

Aqui sou povo sofrido

lá eu serei fazendeiro

terei gado, terei sol

O mar de lá é tão lindo

natureza generosa

que faz nascer sem espinho

o milagre da rosa.

O frio não é muito frio

nem o calor é muito quente

e falam que quem lá vive

é maravilha de gente.
\end{center}

BRANT, F., Nascimento, M. \textbf{Sonho Imigrante}. Disponível em: <http://construtor.aprendebrasil.com.br/up121780001/7910605/Sonho%2520Imigrante.pdf>. Acesso em: 10 dez. 2020. (Adaptado)

\par

A partir de reflexões sobre as leituras do livro \textit{A menina que fez a América}, de Ilka B. Laurito, e do poema \textit{Sonho Imigrante}, de Fernando Brant e de Milton Nascimento, analise as afirmativas abaixo:

\par

I. As pessoas cogitavam deixar as próprias pátrias, a fim de se tornarem cidadãos melhores.

II. Os europeus deixavam seus países esperançosos por uma vida repleta de possibilidades.

III. A busca pelo novo mundo relacionava-se à rejeição do país de origem dos imigrantes.

\par

Está CORRETO o que se afirma em:
\begin{enumerate}
\item I e II, apenas.
\item II, apenas.
\item II e III, apenas.
\item I, II e III.
\end{enumerate}
